% ---- pg-275 （书页 263） ----

\begin{mvtable}
{\zihao{6}\hfill 续表\par}\vspace{2pt}
\begin{tabular}{cccc}
\toprule
\textbf{WordDiff} & \textbf{WordSame} & \textbf{ArabicDiff} & \textbf{ArabicSame} \\
$(\mathbf{x}_{1})$ & $(\mathbf{x}_{2})$ & $(\mathbf{x}_{3})$ & $(\mathbf{x}_{4})$ \\
\midrule
813.5 & 750.5 & 711.0 & 625.0 \\
1 289.5 & 1 140.0 & 904.5 & 784.5 \\
1 096.5 & 1 009.0 & 1 076.0 & 983.0 \\
1 083.0 & 958.0 & 918.0 & 746.5 \\
1 114.0 & 1 046.0 & 1 081.0 & 796.0 \\
708.0 & 669.0 & 657.0 & 572.5 \\
1 201.0 & 925.0 & 1 004.5 & 673.5 \\
\bottomrule
\end{tabular}
\end{mvtable}

{\zihao{6} 资料来源：数据经 J.Carr 许可使用.}

两个数字以单词（“two”，“four”）或者阿拉伯数字（“2”，“4”）形式出现，32 名被测试者需要对此作出一系列迅速的数字判断. 如果两个数字具有相同的奇偶性（都是偶数或者奇数），被测试者回答“相同”；如果它们的奇偶性不同，回答“不同”. 一半被测试者先做阿拉伯数字测试，然后是单词数字测试，另一半被测试者的测试顺序颠倒一下. 对于奇偶性和数字形式的四种组合的每一个，记录每位被测试者的正确反应时间的中值. 设

$X_{1}=$ 单词形式—奇偶性不同的中值反应时间\\
$X_{2}=$ 单词形式—奇偶性相同的中值反应时间\\
$X_{3}=$ 阿拉伯数字形式—奇偶性不同的中值反应时间\\
$X_{4}=$ 阿拉伯数字形式—奇偶性相同的中值反应时间

(a) 用重复测量设计检验处理效应. 令 $\alpha=0.05$.

(b) 对描述数字形式效应、奇偶性效应和它们的交互效应的差异，构造它们的 95\%（联合）置信区间. 解释所得的区间.

(c) 交互作用的缺乏支持数字认知的 M 模型，而交互作用的存在支持数字识别的 C 和 C 模型. 在这个试验中，哪个模型会被接受?

(d) 对于每个被测试者，构造三个对应于数字形式对比、奇偶性对比和交互作用对比的差异成绩. 对于这些数据，多元变量的正态分布是不是一个合理的总体模型? 请解释.

\medskip
\noindent 6.18\quad 朱利可和莫西曼曾对彩龟的大小与形状之间的关系进行过研究（参见文献[12]）. \hyperref[tab:6.9]{表 6.9} 列出了 24 只雄龟与 24 只雌龟的测量值.

\begin{mvtable}
\phantomsection\label{tab:6.9}\mvtabcap{表 6.9\quad 彩龟的龟壳测量值(毫米)}
\begin{tabular}{ccc|ccc}
\toprule
\multicolumn{3}{c}{雌龟} & \multicolumn{3}{c}{雄龟} \\
\cmidrule(lr){1-3}\cmidrule(lr){4-6}
长$(x_{1})$ & 宽$(x_{2})$ & 高$(x_{3})$ & 长$(x_{1})$ & 宽$(x_{2})$ & 高$(x_{3})$ \\
\midrule
98 & 81 & 38 & 93 & 74 & 37 \\
103 & 84 & 38 & 94 & 78 & 35 \\
103 & 86 & 42 & 96 & 80 & 35 \\
105 & 86 & 42 & 101 & 84 & 39 \\
109 & 88 & 44 & 102 & 85 & 38 \\
123 & 92 & 50 & 103 & 81 & 37 \\
123 & 95 & 46 & 104 & 83 & 39 \\
133 & 99 & 51 & 106 & 83 & 39 \\
\bottomrule
\end{tabular}
\end{mvtable}
